\chapter{The cross product : \(\vec{P}\times\vec{Q}\)}

Earlier, we stated that the cross product components of \(\vec{P}\times\vec{Q}\) can be calculated by evaluating the
determinant of the following matrix, which we will arbitrarily call \(A\) here.

\begin{equation*}
    A = \begin{vmatrix}
        \qi & \qj & \quaternion{k}  \\
        P_1      & P_2      & P_3       \\
        Q_1      & Q_2      & Q_3
        \end{vmatrix}
\end{equation*}
\vspace{0.0\baselineskip}

Calculating the determinant of this matrix proceeds as follows;

\begin{equation*}
    \det{(A)} = \qi
    \begin{vmatrix}
        P_2        & P_3  \\
        Q_2        & Q_3
    \end{vmatrix}
    -
    \qj
    \begin{vmatrix}
        P_1        & P_3  \\
        Q_1        & Q_3
    \end{vmatrix}
    +
    \quaternion{k}
    \begin{vmatrix}
        P_1        & P_2  \\
        Q_1        & Q_2
    \end{vmatrix}
\end{equation*}
\vspace{0.0\baselineskip}

Expanding this out gives the following;

\begin{equation*}
    \det{(A)} = \qi (P_2Q_3 - P_3Q_2) - \qj (P_1Q_3 - P_3Q_1) + \quaternion{k}(P_1Q_2 - P_2Q_1)
\end{equation*}
\vspace{0.0\baselineskip}

Which then simplifies one step further to become;

\begin{equation*}
    \det{(A)} = \qi (P_2Q_3 - P_3Q_2) + \qj (P_3Q_1 - P_1Q_3) + \quaternion{k}(P_1Q_2 - P_2Q_1)
\end{equation*}


\section{A very brief history of how the cross product came about.}

If you have ever studied the cross product before and wondered why it is defined in the manner in which it is -- this is
it. That is, it came about in large part thanks to William Rowan Hamilton and his development of quaternions and
quaternion multiplication along with the rules governing it, in the 1840s. We use the expression
``in large part'', because it is widely believed that Hamilton never actually used the term ``cross product'' for that
part of the result which we looked at earlier. Instead, it is reported that he just considered the result of a
quaternion multiplication to be composed of a scalar part and a vector part. Later in the 1800s, Oliver Heaviside and
Josiah Willard Gibbs built on Hamilton's work with quaternions to formalise the mathematical concept of vector. It is
from their work, that the idea of the cross product came about and was subsequently named as such. \\

However, other figures were also involved in work which was relevant to the development of the cross product. In the
1770s, Joseph Louis Lagrange developed formulas which are algebraically equivalent to the cross product. Around the same
time that Hamilton developed the quaternion, Hermann Grassmann also undertook work which helped contribute to the
cross product as we know it today.


\section{What does the cross product actually mean?}

So far in this article, we have talked somewhat at length about the cross product. But what does the cross product
actually mean? To try and help explain what it means, we will consider two different scenarios. The first scenario
involves the multiplication of two separate vectors, while the second scenario involves the multiplication of two
separate quaternions.  \\

\textbullet \; Multiplying two vectors together.  \\

If we calculate the cross product between the two 3-dimensional vectors \(\vec{A}\) and \(\vec{B}\), then the result
will be a third vector which points in a direction that is perpendicular to both \(\vec{A}\) and \(\vec{B}\).  The
magnitude of the resulting vector is given by the following equation\\

\begin{equation}
    |A||B|\sin(\theta)
\end{equation}

\textbullet \; Multiplying two quaternions together.  \\

But what if we multiply two quaternions together, say \(P\) and \(Q\)? In this case, the cross product component of the
result, represents that component of the result which is projected onto the 3-dimensional imaginary subspace of the
4-dimensional quaternion space \(\mathbb{H}\).